
\section{7. Results}

\subsection{7.1 Main Results}

\textbf{Table 7.1: SST-2 results with $S_\text{max} = {{MAX_SPAN}}$, $R = {{RAMP_SIZE}}$, $\lambda = {{LAMBDA}}$.}

\begin{table}[H]
\centering
\begin{tabular}{llllllll}
\toprule
\textbf{Model} & \textbf{Module} & \textbf{Acc} & \textbf{$\Delta$} & \textbf{F1} & \textbf{Learned span} & \textbf{Attn mem} & \textbf{Time} \\
\midrule
DistilBERT & full & 90.71\% & -- & 0.9099 & -- & 128 KB & 1.9 s \\
DistilBERT & **ASA-Span** & **{{D\_ACC}}** & {{D\_DELTA}} & {{D\_F1}} & **{{D\_SPAN}}** & {{D\_MEM}} & {{D\_TIME}} s \\
RoBERTa & full & 78.56\% & -- & 0.8124 & -- & 128 KB & 4.1 s \\
RoBERTa & **ASA-Span** & **{{R\_ACC}}** & {{R\_DELTA}} & {{R\_F1}} & **{{R\_SPAN}}** & {{R\_MEM}} & {{R\_TIME}} s \\
GPT-2 & full & 79.36\% & -- & 0.7978 & -- & 128 KB & 4.2 s \\
GPT-2 & **ASA-Span** & **{{G\_ACC}}** & {{G\_DELTA}} & {{G\_F1}} & **{{G\_SPAN}}** & {{G\_MEM}} & {{G\_TIME}} s \\
\bottomrule
\end{tabular}
\end{table}

\subsection{7.2 Learned Spans}

\textbf{Table 7.2: Span distribution across heads.}

\begin{table}[H]
\centering
\begin{tabular}{llllll}
\toprule
\textbf{Model} & \textbf{Mean span} & \textbf{Min} & \textbf{Max} & \textbf{Std} & \textbf{Reduction from $S_\text{max}$} \\
\midrule
DistilBERT & **{{D\_SPAN}}** & {{D\_SPAN\_MIN}} & {{D\_SPAN\_MAX}} & {{D\_SPAN\_STD}} & {{D\_SPAN\_REDUCTION}}\% \\
RoBERTa & **{{R\_SPAN}}** & {{R\_SPAN\_MIN}} & {{R\_SPAN\_MAX}} & {{R\_SPAN\_STD}} & {{R\_SPAN\_REDUCTION}}\% \\
GPT-2 & **{{G\_SPAN}}** & {{G\_SPAN\_MIN}} & {{G\_SPAN\_MAX}} & {{G\_SPAN\_STD}} & {{G\_SPAN\_REDUCTION}}\% \\
\bottomrule
\end{tabular}
\end{table}

\textbf{Observation.} Learned spans are \textbf{highly non-uniform}. Some heads 
converge to < {{D\_SPAN\_MIN}} (local), others to near $S_\text{max}$ (global).

\subsection{7.3 Per-Layer Span Evolution}

\textbf{Table 7.3: Mean span per layer (DistilBERT).}

\begin{table}[H]
\centering
\begin{tabular}{ll}
\toprule
\textbf{Layer} & \textbf{Mean span} \\
\midrule
0 & {{D\_L0\_SPAN}} \\
1 & {{D\_L1\_SPAN}} \\
2 & {{D\_L2\_SPAN}} \\
3 & {{D\_L3\_SPAN}} \\
4 & {{D\_L4\_SPAN}} \\
5 & {{D\_L5\_SPAN}} \\
\bottomrule
\end{tabular}
\end{table}

\textbf{Observation.} Early layers tend to have \textbf{larger} spans; late layers 
have \textbf{smaller} spans. This matches the observation that early layers 
integrate global context while late layers refine local features.

\subsection{7.4 Efficiency}

\textbf{Table 7.4: Memory and time.}

\begin{table}[H]
\centering
\begin{tabular}{llll}
\toprule
\textbf{Model} & \textbf{Config} & \textbf{Peak VRAM} & \textbf{Inference Time} \\
\midrule
DistilBERT & full & 183 MB & 1.90 s \\
DistilBERT & ASA-Span & {{D\_VRAM}} MB & {{D\_TIME}} s \\
GPT-2 & full & 793 MB & 4.20 s \\
GPT-2 & ASA-Span & {{G\_VRAM}} MB & {{G\_TIME}} s \\
\bottomrule
\end{tabular}
\end{table}

\textbf{Note.} ASA-Span does not reduce inference time in the current implementation 
because the ramp is applied as a \textbf{soft mask} rather than a hard window. 
To realize speedup, a \textbf{hard truncation} post-training is needed.

\subsection{7.5 Ablation: Max Span}

\textbf{Table 7.5: Accuracy vs $S_\text{max}$.}

\begin{table}[H]
\centering
\begin{tabular}{lllll}
\toprule
\textbf{$S_\text{max}$} & \textbf{DistilBERT} & \textbf{RoBERTa} & \textbf{GPT-2} & \textbf{Learned mean} \\
\midrule
32 & {{D\_MAX32}} & {{R\_MAX32}} & {{G\_MAX32}} & {{D\_LEARN32}} \\
64 & {{D\_MAX64}} & {{R\_MAX64}} & {{G\_MAX64}} & {{D\_LEARN64}} \\
128 & {{D\_MAX128}} & {{R\_MAX128}} & {{G\_MAX128}} & {{D\_LEARN128}} \\
256 & {{D\_MAX256}} & {{R\_MAX256}} & {{G\_MAX256}} & {{D\_LEARN256}} \\
\bottomrule
\end{tabular}
\end{table}

\subsection{7.6 Ablation: Ramp Size}

\textbf{Table 7.6: Accuracy vs $R$.}

\begin{table}[H]
\centering
\begin{tabular}{llll}
\toprule
\textbf{$R$} & \textbf{DistilBERT} & \textbf{RoBERTa} & \textbf{GPT-2} \\
\midrule
1 & {{D\_R1}} & {{R\_R1}} & {{G\_R1}} \\
8 & {{D\_R8}} & {{R\_R8}} & {{G\_R8}} \\
16 & {{D\_R16}} & {{R\_R16}} & {{G\_R16}} \\
32 & {{D\_R32}} & {{R\_R32}} & {{G\_R32}} \\
64 & {{D\_R64}} & {{R\_R64}} & {{G\_R64}} \\
\bottomrule
\end{tabular}
\end{table}

\textbf{Observation.} Larger $R$ gives smoother transitions but less aggressive 
pruning.

\subsection{7.7 Ablation: L1 Penalty}

\textbf{Table 7.7: Accuracy and span vs $\lambda$.}

\begin{table}[H]
\centering
\begin{tabular}{llll}
\toprule
\textbf{$\lambda$} & \textbf{DistilBERT acc} & \textbf{Mean span} & \textbf{Reduction} \\
\midrule
0.0 & {{D\_L0}} & {{D\_SPAN0}} & 0\% \\
0.001 & {{D\_L001}} & {{D\_SPAN001}} & {{D\_RED001}}\% \\
0.01 & {{D\_L01}} & {{D\_SPAN01}} & {{D\_RED01}}\% \\
0.1 & {{D\_L1}} & {{D\_SPAN1}} & {{D\_RED1}}\% \\
\bottomrule
\end{tabular}
\end{table}

\textbf{Observation.} $\lambda = 0.01$ is a good balance.

\subsection{7.8 Ablation: LoRA Rank}

\textbf{Table 7.8: Accuracy vs LoRA rank.}

\begin{table}[H]
\centering
\begin{tabular}{llll}
\toprule
\textbf{Rank} & \textbf{DistilBERT} & \textbf{RoBERTa} & \textbf{GPT-2} \\
\midrule
4 & {{D\_LR4}} & {{R\_LR4}} & {{G\_LR4}} \\
8 & {{D\_LR8}} & {{R\_LR8}} & {{G\_LR8}} \\
16 & {{D\_LR16}} & {{R\_LR16}} & {{G\_LR16}} \\
32 & {{D\_LR32}} & {{R\_LR32}} & {{G\_LR32}} \\
\bottomrule
\end{tabular}
\end{table}

\subsection{7.9 Ablation: Recovery Steps}

\textbf{Table 7.9: Accuracy vs recovery steps.}

\begin{table}[H]
\centering
\begin{tabular}{llll}
\toprule
\textbf{Steps} & \textbf{DistilBERT} & \textbf{RoBERTa} & \textbf{GPT-2} \\
\midrule
10 & {{D\_S10}} & {{R\_S10}} & {{G\_S10}} \\
25 & {{D\_S25}} & {{R\_S25}} & {{G\_S25}} \\
50 & {{D\_S50}} & {{R\_S50}} & {{G\_S50}} \\
100 & {{D\_S100}} & {{R\_S100}} & {{G\_S100}} \\
\bottomrule
\end{tabular}
\end{table}

\subsection{7.10 Comparison with Companion Methods}

\textbf{Table 7.10: Comparison.}

\begin{table}[H]
\centering
\begin{tabular}{lllllll}
\toprule
\textbf{Method} & \textbf{DistilBERT} & \textbf{RoBERTa} & \textbf{GPT-2} & \textbf{Type} & \textbf{Params $\Delta$} & \textbf{Attn mem} \\
\midrule
full & 90.71\% & 78.56\% & 79.36\% & Dense & 1.00$\times$ & 1.00$\times$ \\
GQA & {{GQA\_D}} & {{GQA\_R}} & {{GQA\_G}} & KV reduce & 0.95$\times$ & 1.00$\times$ \\
Primal & {{PRIMAL\_D}} & {{PRIMAL\_R}} & {{PRIMAL\_G}} & Low-rank & 1.00$\times$ & **{{PRIMAL\_MEM}}$\times$** \\
ASA-Sparse & {{ASA\_D}} & {{ASA\_R}} & {{ASA\_G}} & Sparse & 1.00$\times$ & 1.00$\times$ \\
**ASA-Span** & **{{D\_ACC}}** & **{{R\_ACC}}** & **{{G\_ACC}}** & Learned span & 1.00$\times$ & 1.00$\times$ \\
\bottomrule
\end{tabular}
\end{table}

\subsection{7.11 Training Loss and Span Evolution}

\textbf{Figure 7.1} shows span evolution over training steps.

\textbf{Observation.} Spans start near $S_\text{max}$ and gradually shrink. The 
shrinkage is faster in early layers, slower in late layers. By step 50, 
convergence is reached.

\subsection{7.12 Summary}

\begin{itemize}
\item \textbf{Accuracy:} drops {{D\_DELTA}}\%--{{G\_DELTA}}\%.
\end{itemize}
\begin{itemize}
\item \textbf{Learned spans:} {{D\_SPAN}}--{{G\_SPAN}} (from max {{MAX\_SPAN}}).
\end{itemize}
\begin{itemize}
\item \textbf{Reduction:} {{D\_SPAN\_REDUCTION}}\%--{{G\_SPAN\_REDUCTION}}\%.
\end{itemize}
\begin{itemize}
\item \textbf{No parameter reduction.}
\end{itemize}
\begin{itemize}
\item \textbf{No speedup in current implementation.}
\end{itemize}
\begin{itemize}
\item \textbf{All models pass quality gate.}
\end{itemize}